\subsection{Training Pipeline}\label{subsec:training-pipeline}

This section describes how we construct supervised training pairs from weekly DeXposure snapshots, define the multi-task learning objective, and optimize DeXposure-FM.

\subsubsection{Data Preparation}\label{subsubsec:data-prep}
DeXposure provides weekly snapshots of the inter-protocol exposure network. For each discrete time index $\tau$, we construct a weighted directed graph
$G_{\tau}=(\mathcal{P}_{\tau},\mathcal{E}_{\tau})$ with node weights $w_{\tau}(p)$ (protocol TVL) and edge weights $w_{\tau}(e_{pq})$ (exposure from $p$ to $q$).
Training examples are formed by an \emph{anchor} snapshot and a forecast horizon:
\begin{equation}
G_{\tau} \;\longrightarrow\; G_{\tau+h},
\qquad
h\in\{1,4,8,12\},
\label{eq:train-pairs}
\end{equation}
so that the model predicts the future graph state at horizon $h$ conditional on the current snapshot.

For each protocol $p\in\mathcal{P}_{\tau}$ we provide tabular node features $x^{\text{tab}}_{p,\tau}$ (including log-scaled TVL, token counts, concentration measures, and sector/category indicators) together with the graph structure $G_{\tau}$.

Table~\ref{tab:data-stats} summarizes the dataset characteristics. The network exhibits high edge persistence (mean overlap ratio 98.5\%), reflecting the stable nature of DeFi protocol interconnections.

\begin{table}[t]
\centering
\small
\begin{tabular}{lr}
\hline
\textbf{Statistic} & \textbf{Value} \\
\hline
Total snapshots & 283 weeks \\
Time range & 2020-03-23 $\sim$ 2025-08-18 \\
Mean nodes per week & 5,676 \\
Mean edges per week & 30,424 \\
Mean edge overlap ratio & 98.5\% \\
Protocol categories & $\sim$ \\
\hline
\end{tabular}
\caption{DeXposure dataset statistics.}
\label{tab:data-stats}
\end{table}

\textbf{Time-based expanding window splits.}
We adopt an expanding-window walk-forward evaluation protocol to prevent look-ahead bias: all validation and test targets occur strictly after the corresponding training window, and the training window expands over time while validation and test windows remain fixed.

\textbf{Negative sampling.}
Edge existence is highly imbalanced.  For each horizon $h$, we construct a labeled edge set $\mathcal{S}_{\tau+h}$ consisting of all positive edges in $\mathcal{E}_{\tau+h}$ plus uniformly sampled negative edges from the complement, using a negative-to-positive ratio of $5{:}1$.
This sampled set is used for the edge existence loss, while the edge weight loss is computed only on positive edges.

\subsubsection{Loss Functions}\label{subsubsec:loss}
DeXposure-FM is trained with a multi-task objective combining (1) edge existence classification, (2) edge weight regression, and (3) node-level TVL-change prediction.

\textbf{Edge existence (link prediction).}
For each candidate pair $(p,q)\in\mathcal{S}_{\tau+h}$ we predict $\widehat{y}^{\text{exist}}_{pq}\in(0,1)$ and minimize a weighted binary cross-entropy loss:
\begin{equation}
\mathcal{L}_{\text{exist}}
=
\sum_{(p,q)\in\mathcal{S}_{\tau+h}}
\mathrm{BCE}\!\left(\widehat{y}^{\text{exist}}_{pq},\,y^{\text{exist}}_{pq};\,w_{\text{pos}}\right),
\label{eq:loss-exist}
\end{equation}
where $w_{\text{pos}}$ upweights positives to match the negative sampling ratio. Here, binary cross-entropy is defined as below. 

\begin{definition}[binary cross-entropy (BCE)]
For a binary label $y\in\{0,1\}$ and a predicted probability $\hat{y}\in(0,1)$, the binary cross-entropy loss is
\begin{equation}
\mathrm{BCE}(\hat{y},y)
= -\Big(y\log(\hat{y}) + (1-y)\log\big(1-\hat{y}\big)\Big).
\label{eq:bce}
\end{equation}
Equivalently, if $\hat{y}=\sigma(z)$ is produced from a logit $z\in\mathbb{R}$ via the sigmoid $\sigma(z)=1/(1+e^{-z})$, then
\begin{equation}
\mathrm{BCE}(\sigma(z),y)
= \log\!\big(1+e^{z}\big) - yz,
\label{eq:bce-logit}
\end{equation}
which is numerically stable and is often referred to as the logistic loss.
\end{definition}

\textbf{Edge weights (residual learning).}
For positive edges $(p,q)\in\mathcal{E}_{\tau+h}$ we predict a residual in log space relative to the previous observed week. Let $\tilde{w}_{\tau}(e_{pq})=\log(1+w_{\tau}(e_{pq}))$ and $\tilde{w}_{\tau+h}(e_{pq})=\log(1+w_{\tau+h}(e_{pq}))$. The link head outputs $\widehat{r}_{pq,\tau,h}$ and we apply a robust Smooth~L1 loss:
\begin{equation}
\mathcal{L}_{\text{weight}}
=
\sum_{(p,q)\in\mathcal{E}_{\tau+h}}
\mathrm{SmoothL1}\!\Big(
\widehat{r}_{pq,\tau,h}
-
\big(\tilde{w}_{\tau+h}(e_{pq})-\tilde{w}_{\tau}(e_{pq})\big)
\Big).
\label{eq:loss-weight}
\end{equation}
Equivalently, we reconstruct $\widehat{\tilde{w}}_{\tau+h}(e_{pq})=\tilde{w}_{\tau}(e_{pq})+\widehat{r}_{pq,\tau,h}$ and minimize $\mathrm{SmoothL1}(\widehat{\tilde{w}}_{\tau+h}(e_{pq})-\tilde{w}_{\tau+h}(e_{pq}))$. We set $\tilde{w}_{\tau}(e_{pq})=0$ when the edge is absent at time $\tau$. The $\log(1+\cdot)$ transform stabilizes training under the heavy-tailed exposure distribution. Here, smooth L1 (Huber) loss is defined as follows.
\begin{definition}[smooth L1 (Huber) loss]
For a scalar prediction $\hat{u}\in\mathbb{R}$ and target $u\in\mathbb{R}$, let the residual be $r=\hat{u}-u$.
The Smooth~L1 (Huber) loss with threshold $\delta>0$ is defined as
\begin{equation}
\mathrm{SmoothL1}_\delta(r)=
\begin{cases}
\frac{1}{2}\frac{r^2}{\delta}, & \text{if } |r|<\delta,\\[6pt]
|r|-\frac{1}{2}\delta, & \text{otherwise}.
\end{cases}
\label{eq:smoothl1}
\end{equation}
Equivalently, written directly in terms of $(\hat{u},u)$:
\begin{equation}
\mathrm{SmoothL1}_\delta(\hat{u},u)=
\begin{cases}
\frac{1}{2}\frac{(\hat{u}-u)^2}{\delta}, & \text{if } |\hat{u}-u|<\delta,\\[6pt]
|\hat{u}-u|-\frac{1}{2}\delta, & \text{otherwise}.
\end{cases}
\end{equation}
\end{definition}
Smooth~L1 behaves like an $\ell_2$ loss near zero (encouraging small residuals) and like an $\ell_1$ loss for large residuals (reducing sensitivity to heavy tails and outliers). A common choice is $\delta=1$, when $\mathrm{SmoothL1}_\delta$ is briefly written as $\mathrm{SmoothL1}$.

\textbf{Node-level TVL dynamics.}
We also predict the log-change in protocol TVL,
$\Delta_{p,\tau+h}=\log(1+w_{\tau+h}(p))-\log(1+w_{\tau}(p))$,
using Smooth~L1:
\begin{equation}
\mathcal{L}_{\text{node}}
=
\sum_{p\in\mathcal{P}_{\tau}\cap \mathcal{P}_{\tau+h}}
\mathrm{SmoothL1}\!\left(\widehat{\Delta}_{p,\tau+h}-\Delta_{p,\tau+h}\right),
\label{eq:loss-node}
\end{equation}
where $\Delta_{p,\tau+h} = \log(1+w_{\tau+h}(p)) - \log(1+w_{\tau}(p))$ is the ground-truth log-change in TVL.


\textbf{Multi-task objective.}
The total loss is a weighted sum:
\begin{equation}
\mathcal{L}
=
\lambda_{\text{exist}}\mathcal{L}_{\text{exist}}
+
\lambda_{\text{weight}}\mathcal{L}_{\text{weight}}
+
\lambda_{\text{node}}\mathcal{L}_{\text{node}},
\label{eq:loss-total}
\end{equation}
with $\lambda_{\text{exist}}=2.0$, $\lambda_{\text{weight}}=0.5$, and $\lambda_{\text{node}}=20.0$. These weights are calibrated to balance the gradient contributions across tasks despite their different loss magnitudes: edge existence loss $\mathcal{L}_{\mathrm{exist}} \approx 0.22$, edge weight loss $\mathcal{L}_{\mathrm{weight}} \approx 2.4$, and node loss $\mathcal{L}_{\mathrm{node}} \approx 0.05$ on average. They also reflect the primary importance of link prediction while using weight and node supervision to encourage economically meaningful representations. 

\subsubsection{Optimization}\label{subsubsec:optim}

\textbf{Initialization.} We adopt the open-source weights of GraphPFN as initialization of the encoder. 

\textbf{Optimizer.}
We use Adam \citep{KingmaBa2015Adam} with hyperparameters $\beta_1=0.9$, $\beta_2=0.999$ and learning rate $\eta = 5 \times 10^{-4}$.

\textbf{Training details.}
We train for up to 20 epochs with early stopping (patience = 3) based on validation AUPRC. Gradient clipping ($\|\nabla\|_2 \le 1.0$) is applied for stability.


\textbf{Optional sharpness-aware training.}
Sharpness-Aware Minimization (SAM) is a promising extension to improve robustness under non-stationarity by favoring flatter optima. However, our current released training pipeline and the experiments reported in this paper use Adam; we leave SAM-based training as future work.

\begin{table}[t]
\centering
\small
\setlength{\tabcolsep}{6pt}
\renewcommand{\arraystretch}{1.05}
\begin{tabular}{ll}
\hline
\textbf{Hyperparameter} & \textbf{Value} \\
\hline
Data granularity & Weekly snapshots \\
Forecast horizons $h$ & $\{1, 4, 8, 12\}$ weeks \\
Negative sampling ratio & 5:1 (neg:pos) \\
Optimizer & Adam ($\beta_1=0.9$, $\beta_2=0.999$) \\
Learning rate (heads) & $5 \times 10^{-4}$ \\
Learning rate (backbone) & $5 \times 10^{-5}$ (fine-tune only) \\
Training epochs & 20 (early stopping patience=3) \\
Gradient clipping & $\|\nabla\|_2 \le 1.0$ \\
\hline
\multicolumn{2}{l}{\textit{Loss weights:}} \\
$\lambda_{\mathrm{exist}}$ (edge existence) & $2.0$ \\
$\lambda_{\mathrm{weight}}$ (edge weight) & $0.5$ \\
$\lambda_{\mathrm{node}}$ (node prediction) & $20.0$ \\
\hline
\end{tabular}
\caption{Training hyperparameters for DeXposure-FM experiments.}
\label{tab:training-hparams}
\end{table}
